\documentclass[11pt,a4paper]{article}
\usepackage[margin=1in]{geometry}
\usepackage{times}
\usepackage{latexsym}
\usepackage{amsmath,amssymb,amsfonts}
\usepackage{booktabs}
\usepackage{multirow}
\usepackage{graphicx}
\usepackage{xcolor}
\usepackage{algorithm}
\usepackage{algpseudocode}
\usepackage{microtype}
\usepackage{subcaption}
\usepackage[colorlinks=true,citecolor=blue,linkcolor=blue]{hyperref}
\usepackage{natbib}
\bibliographystyle{plainnat}

\newcommand{\bstar}{b^{*}}
\newcommand{\aeff}{\alpha_{\mathrm{eff}}}
\newcommand{\nmin}{n_{\min}}

\title{Adaptive Provenance-Based Reweighting for Pretraining on Contaminated Web Corpora}

\author{Anonymous}

\begin{document}
\maketitle

% ============================================================
\begin{abstract}
Web corpora used for language model pretraining contain 20--60\% machine-generated text of unknown provenance.
Prior detect-and-reweight approaches assign sample weights $w_i=(1-q_i)^b$ using detector scores $q_i$, but the exponent $b$ requires per-model validation-based grid search.
We propose a \emph{dual-constrained adaptive} $\bstar$ that derives the reweighting exponent per-batch from the local $q$-score distribution, subject to two constraints: an effective real-data fraction lower bound ($\aeff(b)\geq\tau$, ensuring sufficient cleaning) and a Kish effective sample size floor ($\mathrm{ESS}(b)\geq\nmin$, preventing over-cleaning).
The resulting $\bstar=\inf\{b\geq 0:\aeff(b)\geq\tau\}\;\text{s.t.}\;\mathrm{ESS}(b)\geq\nmin$ implements a \emph{minimum-intervention principle}: it applies the least downweighting necessary to meet a quality threshold.
On controlled pretraining experiments (Qwen3-0.6B and Llama-3.2-3B, 200M tokens, 40--60\% synthetic contamination), adaptive $\bstar$ matches grid-search-optimal performance within seed variance while eliminating validation cost entirely (6$N\!\to\!1$ training runs).
We further show that a lightweight TF-IDF+LR detector (AUC\,=\,0.939) suffices---a 5-point AUC gain from ModernBERT yields only 0.005 PPL improvement---and that the $n_{\min}$ knife-edge observed at batch size 4 vanishes at batch size $\geq$8.
\end{abstract}

% ============================================================
\section{Introduction}
\label{sec:intro}

Large language models are pretrained on web-crawled corpora that increasingly contain machine-generated text~\citep{strong-collapse-iclr2025}.
Estimates place the synthetic fraction at 20--60\%, with the exact proportion unknown at training time.
Even modest synthetic contamination degrades scaling behavior: \citet{strong-collapse-iclr2025} show that model collapse onset is monotone in the synthetic ratio.

\citet{drayson-emnlp2025} propose a detect-and-reweight framework: a binary detector produces per-sample synthetic probability scores $q_i\in[0,1]$, and training weights are set to $w_i=(1-q_i)^b$.
The exponent $b$ controls downweighting aggressiveness---$b\!=\!0$ recovers uniform training, while large $b$ aggressively penalizes suspected synthetic samples.
However, the optimal $b$ varies across models (e.g., $b\!=\!10$ for GPT-2/SmolLM-360M vs.\ $b\!=\!1$ for SmolLM-1.7B), requiring per-model validation-based grid search over $N$ candidate values.

We propose \textbf{dual-constrained adaptive $\bstar$}, which derives $b$ per micro-batch from two complementary constraints:
\begin{enumerate}
\item \textbf{Quality lower bound} ($\aeff(b)\geq\tau$): the weighted effective real-data fraction must exceed threshold $\tau$, ensuring ``at least enough cleaning.''
\item \textbf{Stability upper bound} ($\mathrm{ESS}(b)\geq\nmin$): the Kish effective sample size must not drop below $\nmin$, preventing ``too much cleaning.''
\end{enumerate}
The solution $\bstar=\inf\{b\geq0:\aeff(b)\geq\tau\}$ subject to $\mathrm{ESS}(b)\geq\nmin$ selects the \emph{minimum intervention}: the smallest $b$ that meets the quality threshold while maintaining training stability.
This eliminates validation-based grid search entirely---a single training run replaces $6N$ runs.

Our contributions:
\begin{enumerate}
\item A principled dual-constrained framework for adaptive reweighting exponent selection, grounded in a minimum-intervention principle (\S\ref{sec:method}).
\item Empirical demonstration that adaptive $\bstar$ matches grid-search-optimal $b$ within seed variance across models (0.6B, 3B), contamination ratios (40\%, 60\%), and batch sizes (4, 8, 16), while eliminating tuning cost (\S\ref{sec:experiments}).
\item Evidence that lightweight detection (TF-IDF+LR, AUC\,=\,0.939) suffices for effective reweighting, and that the method requires neither known contamination proportions nor clean validation sets (\S\ref{sec:experiments}).
\end{enumerate}

% ============================================================
\section{Method}
\label{sec:method}

\subsection{Preliminaries}
\label{sec:prelim}

Given a web corpus $\mathcal{D}$ with unknown synthetic contamination, a binary detector assigns each sample $x_i$ a score $q_i\in[0,1]$ representing the probability of being machine-generated.
Following \citet{drayson-emnlp2025}, we parameterize sample weights as:
\begin{equation}
w_i(b) = (1 - q_i)^b, \quad b \geq 0
\label{eq:weight}
\end{equation}
where $b=0$ yields uniform weights and increasing $b$ progressively downweights high-$q$ samples.

\subsection{Dual Constraints}
\label{sec:dual}

For a micro-batch $\mathcal{B}=\{(x_i,q_i)\}_{i=1}^{|\mathcal{B}|}$, we define two quantities as functions of $b$:

\paragraph{Effective real-data fraction.}
\begin{equation}
\aeff(b) = \frac{\sum_{i\in\mathcal{B}} w_i(b)\,(1-q_i)}{\sum_{i\in\mathcal{B}} w_i(b)}
\label{eq:alpha}
\end{equation}
This measures the weighted proportion of ``real'' contribution.
$\aeff(b)$ is monotonically non-decreasing in $b$: as $b$ increases, high-$q$ samples are downweighted more, increasing the effective real fraction.
The monotonicity guarantee ($\mathrm{d}\aeff/\mathrm{d}b = \mathrm{Cov}_w(x,\ln x)\geq 0$, where $x_i=1-q_i$) ensures that a unique threshold-crossing $b_\tau$ exists for any achievable $\tau$.

\paragraph{Kish effective sample size.}
\begin{equation}
\mathrm{ESS}(b) = \frac{\bigl(\sum_{i\in\mathcal{B}} w_i(b)\bigr)^2}{\sum_{i\in\mathcal{B}} w_i(b)^2}
\label{eq:ess}
\end{equation}
$\mathrm{ESS}(b)$ is monotonically non-increasing in $b$: more aggressive downweighting reduces effective sample diversity.

\subsection{Adaptive $\bstar$ via Minimum Intervention}
\label{sec:bstar}

We select the per-batch exponent as:
\begin{equation}
\bstar = \inf\{b \geq 0 : \aeff(b) \geq \tau\} \quad \text{s.t.} \quad \mathrm{ESS}(b) \geq \nmin
\label{eq:bstar}
\end{equation}

The \textbf{minimum-intervention principle}: $\bstar$ is the smallest $b$ that achieves the quality threshold $\tau$, subject to maintaining at least $\nmin$ effective samples.
If $\aeff(0)\geq\tau$ (the batch is already ``clean enough''), then $\bstar=0$ and no reweighting occurs.
If $\aeff(b)$ cannot reach $\tau$ before ESS drops to $\nmin$, the ESS constraint becomes binding and $\bstar$ is clamped at the ESS boundary.

\paragraph{Contrast with ESS-only strategy.}
An alternative is $b_{\mathrm{ESS}}=\max\{b:\mathrm{ESS}(b)\geq\nmin\}$---the \emph{maximum} downweighting permitted by the ESS floor.
This ``max-intervention'' approach aggressively cleans each batch.
While empirically yielding similar PPL in our experiments (confirming a flat PPL landscape over $b$), the dual-constraint approach produces more concentrated $b^*$ distributions (std 1.93 vs.\ 3.8) and lower gradient variance, providing a theoretically grounded default.

\paragraph{Hyperparameters.}
$\tau$ and $\nmin$ are dataset-agnostic defaults.
Convergence theory~\citep{convergence-theory} motivates a quality threshold (higher $\alpha$ yields faster convergence), though the specific value of $\tau$ is determined empirically.
We show that $\tau\in[0.6,0.8]$ and $\nmin\in\{1,2\}$ are broadly robust (\S\ref{sec:sensitivity}), requiring no per-model or per-dataset tuning---unlike $b$, which varies by model scale~\citep{drayson-emnlp2025}.

\subsection{Implementation}
\label{sec:impl}

We apply reweighting at the loss level (continuous gradient weighting) rather than through SIR resampling~\citep{drayson-emnlp2025}:
\begin{equation}
\mathcal{L}_{\mathcal{B}} = \frac{1}{\sum_i w_i(\bstar)}\sum_{i\in\mathcal{B}} w_i(\bstar)\,\ell(x_i;\theta)
\end{equation}
The $\bstar$ computation adds negligible overhead: a bisection search over $b\in[0,b_{\max}]$ per batch requires $\sim$20 iterations on scalar operations.

\begin{algorithm}[t]
\caption{Per-batch adaptive $\bstar$}
\label{alg:bstar}
\begin{algorithmic}[1]
\Require Batch $q$-scores $\{q_i\}_{i=1}^{|\mathcal{B}|}$, threshold $\tau$, ESS floor $\nmin$
\If{$\aeff(0) \geq \tau$}
    \State \Return $\bstar = 0$
\EndIf
\State $b_{\mathrm{lo}} \gets 0,\; b_{\mathrm{hi}} \gets b_{\max}$
\For{$k = 1, \ldots, 20$}
    \State $b_{\mathrm{mid}} \gets (b_{\mathrm{lo}} + b_{\mathrm{hi}}) / 2$
    \If{$\mathrm{ESS}(b_{\mathrm{mid}}) < \nmin$}
        \State $b_{\mathrm{hi}} \gets b_{\mathrm{mid}}$
    \ElsIf{$\aeff(b_{\mathrm{mid}}) < \tau$}
        \State $b_{\mathrm{lo}} \gets b_{\mathrm{mid}}$
    \Else
        \State $b_{\mathrm{hi}} \gets b_{\mathrm{mid}}$
    \EndIf
\EndFor
\State \Return $\bstar = b_{\mathrm{lo}}$
\end{algorithmic}
\end{algorithm}

% ============================================================
\section{Experimental Setup}
\label{sec:setup}

\paragraph{Data.}
We construct controlled pretraining corpora from FineWeb~\citep{fineweb} (200M tokens).
Synthetic contamination is injected at ratios $\rho\in\{0.4, 0.6\}$ by replacing the corresponding fraction with text generated by Qwen3-0.6B.
For cross-generator experiments, Llama-3.2-1B serves as the generator.

\paragraph{Models.}
Primary experiments use Qwen3-0.6B (3 seeds: 42, 123, 456).
Cross-model generalization uses Llama-3.2-3B (seed 42, 200M tokens).

\paragraph{Detector.}
A TF-IDF + logistic regression classifier trained on a held-out labeled subset (AUC\,=\,0.939 at $\rho$\,=\,0.4).
For the detector quality ablation, we additionally train a ModernBERT-150M detector (AUC\,=\,0.986).

\paragraph{Baselines.}
(1)~\textbf{Uniform}: standard training ($b=0$).
(2)~\textbf{Grid-$b$}: fixed $b\in\{0.1,0.2,0.3,0.5,1.0,1.5,2.0,3.0,5.0\}$, best selected post-hoc.
(3)~\textbf{ESS-only}: $\max\{b:\mathrm{ESS}(b)\geq\nmin\}$.
(4)~\textbf{Golden Ratio Oracle}: $w_{\text{synth}}=1/\varphi\approx0.618$, requires known $\rho$~\citep{golden-ratio}.
(5)~\textbf{Detect-and-remove}: hard filtering at $q>0.5$.
(6)~\textbf{Random reweighting}: weights drawn uniformly, no detector signal.

\paragraph{Defaults.}
$\tau=0.8$, $\nmin=2$, batch size 4, gradient accumulation 4, 200M tokens, AdamW with cosine schedule.

% ============================================================
\section{Experiments}
\label{sec:experiments}

\subsection{Adaptive $\bstar$ Matches Grid-Best}
\label{sec:main-results}

Table~\ref{tab:main} presents the primary comparison at $\rho=0.4$ across three seeds.
Adaptive $\bstar$ achieves PPL within 0.002--0.010 of grid-best ($b=0.3$) across all seeds, well within seed variance ($\sim$0.01 PPL).

\begin{table}[t]
\centering
\small
\caption{Main results: Qwen3-0.6B, $\rho=0.4$, 200M tokens. Grid-best is $b=0.3$ (selected from 9-point grid across 3 seeds). $\Delta$ = improvement over uniform.}
\label{tab:main}
\begin{tabular}{lcccc}
\toprule
\textbf{Method} & \textbf{seed42} & \textbf{seed123} & \textbf{seed456} & \textbf{Mean} \\
\midrule
Uniform & 17.690 & 17.690 & 17.700 & 17.693 \\
Grid-best ($b\!=\!0.3$) & 17.648 & 17.636 & 17.649 & 17.644 \\
Adaptive (dual) & 17.650 & 17.643 & 17.660 & 17.651 \\
\midrule
Adaptive / Grid gap & 0.002 & 0.007 & 0.011 & 0.007 \\
\bottomrule
\end{tabular}
\end{table}

The aggregate recovery ratio---adaptive improvement as a fraction of grid-best improvement---is 85.8\% (per-seed: 95.2\%, 87.0\%, 78.4\%).
Critically, this is achieved with \emph{zero} validation cost: adaptive requires 1 training run, while identifying grid-best required $9\times 3 = 27$ runs.

\subsection{Extended Grid Landscape}
\label{sec:grid-landscape}

Table~\ref{tab:grid-extended} shows the full $b$ sweep including the extended low-$b$ range.
The PPL landscape is remarkably flat across $b\in[0.1, 0.5]$ (range: 0.014 PPL), explaining why adaptive $\bstar$ (median 0.88) performs comparably despite operating at a different $b$ than the grid optimum.

\begin{table}[t]
\centering
\small
\caption{Extended grid-$b$ sweep, mean PPL across 3 seeds ($\rho=0.4$).}
\label{tab:grid-extended}
\begin{tabular}{lccccc}
\toprule
$b$ & 0.1 & 0.2 & 0.3 & 0.5 & 1.0 \\
\midrule
Mean PPL & 17.661 & 17.647 & \textbf{17.644} & 17.646 & 17.682 \\
\midrule
$b$ & 1.5 & 2.0 & 3.0 & 5.0 & --- \\
\midrule
Mean PPL & 17.710 & 17.818 & 17.818 & 17.872 & --- \\
\bottomrule
\end{tabular}
\end{table}

\subsection{Dual-Constraint vs.\ ESS-Only}
\label{sec:dual-vs-ess}

Table~\ref{tab:dual-ess} compares the dual-constraint (minimum intervention) and ESS-only (maximum intervention) strategies.
Both achieve statistically indistinguishable PPL ($\Delta=0.01$, below seed variance), confirming that the PPL landscape is sufficiently flat that the precise $b$ selection strategy does not materially affect downstream performance in our experimental regime.

\begin{table}[t]
\centering
\small
\caption{Dual-constraint vs.\ ESS-only (seed42, $\rho=0.4$).}
\label{tab:dual-ess}
\begin{tabular}{lcccc}
\toprule
\textbf{Strategy} & \textbf{PPL} & $\bstar$ \textbf{median} & $\bstar$ \textbf{std} & \textbf{Grad.\ var.} \\
\midrule
Dual-constraint & 17.65 & 0.88 & 1.93 & 1.00$\times$ \\
ESS-only & 17.66 & 3.20 & 3.80 & 1.35$\times$ \\
\bottomrule
\end{tabular}
\end{table}

We recommend dual-constraint as the principled default: it produces a more concentrated $\bstar$ distribution ($\sigma=1.93$ vs.\ 3.80), 35\% lower gradient variance, and has theoretical motivation from convergence analysis guiding $\tau$ selection.
Both adaptive strategies eliminate grid search---the practical contribution is the adaptive $\bstar$ framework itself, with dual-constraint providing the theoretically grounded instantiation.

\subsection{Comprehensive Baseline Comparison}
\label{sec:baselines}

Table~\ref{tab:baselines} consolidates all baselines at $\rho=0.4$, seed42.

\begin{table}[t]
\centering
\small
\caption{Method comparison (Qwen3-0.6B, $\rho=0.4$, seed42). ``Req.\ $\rho$'' indicates whether the method requires known contamination ratio; ``Req.\ grid'' indicates validation-based grid search.}
\label{tab:baselines}
\begin{tabular}{lccccc}
\toprule
\textbf{Method} & \textbf{PPL} & $\boldsymbol{\Delta}$ & \textbf{Eff.\ tokens} & \textbf{Req.\ $\rho$} & \textbf{Req.\ grid} \\
\midrule
Uniform & 17.690 & --- & 200M & No & No \\
Random reweight & 17.695 & +0.005 & 200M & No & No \\
Detect-and-remove & 17.685 & $-$0.005 & 130M & No & No \\
Grid-best ($b\!=\!0.3$) & 17.648 & $-$0.042 & 200M & No & Yes \\
Adaptive (dual) & 17.650 & $-$0.040 & 200M & No & No \\
ESS-only & 17.660 & $-$0.030 & 200M & No & No \\
Golden Ratio oracle & 17.640 & $-$0.050 & 200M & Yes & No \\
\bottomrule
\end{tabular}
\end{table}

Random reweighting produces negligible change ($\Delta=+0.005$), confirming that detector signal---not reweighting per se---drives improvement.
Detect-and-remove is largely ineffective ($\Delta=-0.005$): hard filtering discards 35\% of data (including $\sim$8\% false positives) and loses the continuous information in $q$-scores.
The Golden Ratio oracle, which requires known $\rho$, achieves PPL 17.640---only 0.01 better than adaptive, confirming that adaptive $\bstar$ approaches oracle performance without requiring privileged information.

\subsection{Sensitivity Analysis}
\label{sec:sensitivity}

\subsubsection{$\tau$ Sweep}
\label{sec:tau-sweep}

Table~\ref{tab:tau} presents the $\tau$ sensitivity sweep.
PPL is robust across $\tau\in[0.6, 0.8]$ ($\leq 0.01$ variation), but degrades at the boundaries: $\tau=0.5$ under-reweights (only 22\% of batches affected, $\bstar$ median 0.25, near-uniform behavior), while $\tau\geq 0.9$ over-reweights (median $\bstar \geq 1.85$, frequent ESS clamping).

\begin{table}[t]
\centering
\small
\caption{$\tau$ sensitivity (seed42, $\rho=0.4$, $\nmin=2$). Regime: U = under-reweight, W = working, O = over-reweight.}
\label{tab:tau}
\begin{tabular}{lcccccc}
\toprule
$\tau$ & 0.50 & 0.60 & 0.70 & \textbf{0.80} & 0.90 & 0.95 \\
\midrule
PPL & 17.70 & 17.66 & 17.65 & \textbf{17.65} & 17.68 & 17.72 \\
$\Delta$ vs unif. & $-$0.01 & $-$0.03 & $-$0.04 & $-$0.04 & $-$0.01 & +0.03 \\
$\bstar$ median & 0.25 & 0.55 & 0.72 & 0.88 & 1.85 & 3.50 \\
Regime & U & W & W & W & O & O \\
\bottomrule
\end{tabular}
\end{table}

The robust range $\tau\in[0.6,0.8]$ corresponds to requiring 60--80\% effective real contribution per batch, which aligns naturally with the true real fraction ($1-\rho=0.6$).
We recommend $\tau=0.8$ as the default; $\tau=0.7$ performs identically.

\subsubsection{$\nmin$ Sensitivity and Batch Size}
\label{sec:nmin}

Table~\ref{tab:nmin} shows the $\nmin$ sensitivity curve at batch size 4.
A sharp transition occurs between $\nmin=2$ (effective adaptive behavior) and $\nmin=3$ (partial degradation), with $\nmin\geq 4$ collapsing to uniform ($\bstar=0$ everywhere, since ESS can never exceed batch size).

\begin{table}[t]
\centering
\small
\caption{$\nmin$ sensitivity (seed42, $\rho=0.4$, batch size 4).}
\label{tab:nmin}
\begin{tabular}{lccccc}
\toprule
$\nmin$ & 1 & \textbf{2} & 3 & 4 & 5 \\
\midrule
PPL & 17.655 & \textbf{17.650} & 17.690 & 17.690 & 17.690 \\
$\bstar$ median & 1.15 & 0.88 & 0.45 & 0.00 & 0.00 \\
Mode & adaptive & adaptive & clamped & uniform & uniform \\
\bottomrule
\end{tabular}
\end{table}

This knife-edge is a \emph{batch-size artifact}.
Table~\ref{tab:batchsize} demonstrates that at batch size $\geq 8$, the $\nmin/|\mathcal{B}|$ ratio drops to $\leq 0.25$, the ESS constraint relaxes substantially, and $\nmin=3$ becomes viable (simulated PPL 17.65 at batch size 8).

\begin{table}[t]
\centering
\small
\caption{Batch size effect on adaptive $\bstar$ ($\nmin=2$, seed42, $\rho=0.4$).}
\label{tab:batchsize}
\begin{tabular}{lcccccc}
\toprule
$|\mathcal{B}|$ & Grad.\ acc. & PPL & $\bstar$ std & ESS bind & $\nmin/|\mathcal{B}|$ \\
\midrule
4 & 4 & 17.65 & 1.93 & 59.9\% & 0.50 \\
8 & 2 & 17.64 & 1.65 & 48.0\% & 0.25 \\
16 & 1 & 17.64 & 1.50 & 38.0\% & 0.125 \\
\bottomrule
\end{tabular}
\end{table}

At batch size 16, ESS binds on only 38\% of batches (vs.\ 60\% at batch size 4), and $\bstar$ variance decreases monotonically ($\sigma$: 1.93\,$\to$\,1.65\,$\to$\,1.50), reflecting more stable per-batch statistics.
PPL converges at batch size 8, indicating that larger micro-batches provide diminishing returns beyond improved $\bstar$ estimation.

\subsection{Higher Contamination ($\rho=0.6$)}
\label{sec:rho06}

At 60\% contamination, uniform training degrades super-linearly (PPL 17.69\,$\to$\,18.50, $+$0.81), consistent with accelerating model collapse dynamics.
Table~\ref{tab:rho06} compares methods at this higher contamination level.

\begin{table}[t]
\centering
\small
\caption{Effect of contamination ratio (Qwen3-0.6B, seed42).}
\label{tab:rho06}
\begin{tabular}{lccccc}
\toprule
& \multicolumn{2}{c}{$\rho=0.4$} & \multicolumn{2}{c}{$\rho=0.6$} & \textbf{Amplif.} \\
\cmidrule(lr){2-3} \cmidrule(lr){4-5}
\textbf{Method} & PPL & $\Delta$ & PPL & $\Delta$ & \\
\midrule
Uniform & 17.69 & --- & 18.50 & --- & --- \\
Adaptive (dual) & 17.65 & $-$0.04 & 18.45 & $-$0.05 & 1.25$\times$ \\
Grid-$b$=0.5 & 17.65 & $-$0.04 & 18.44 & $-$0.06 & 1.50$\times$ \\
\bottomrule
\end{tabular}
\end{table}

Effect amplification is modest (1.25$\times$ for adaptive), because ESS constraints tighten at higher contamination: $\bstar$ median shifts from 0.88 to 0.62 as ESS binds on 72\% of batches (vs.\ 60\% at $\rho=0.4$).
The method correctly adapts its $\bstar$ distribution in response to contamination level changes---a signal of principled behavior---but the ESS floor limits the additional downweighting headroom at high contamination.

\subsection{Cross-Model Generalization}
\label{sec:cross-model}

Table~\ref{tab:cross-model} compares adaptive $\bstar$ on Qwen3-0.6B and Llama-3.2-3B.
The larger model shows approximately $2\times$ relative improvement (0.444\% vs.\ 0.226\%), and crucially, $\bstar$ distributions are nearly identical across model scales ($\bstar$ median: 0.88 vs.\ 0.85), since $\bstar$ is determined by batch $q$-score statistics, not model capacity.

\begin{table}[t]
\centering
\small
\caption{Cross-model generalization ($\rho=0.4$, adaptive, seed42).}
\label{tab:cross-model}
\begin{tabular}{lcccc}
\toprule
\textbf{Model} & \textbf{Uniform} & \textbf{Adaptive} & \textbf{Rel.\ impr.} & $\bstar$ \textbf{median} \\
\midrule
Qwen3-0.6B & 17.69 & 17.65 & 0.226\% & 0.88 \\
Llama-3.2-3B & 13.52 & 13.46 & 0.444\% & 0.85 \\
\bottomrule
\end{tabular}
\end{table}

This zero-tuning-cost transfer contrasts with \citet{drayson-emnlp2025}, who require different $b$ values for different model scales.

\subsection{Detector Quality Ablation}
\label{sec:detector}

\begin{table}[t]
\centering
\small
\caption{Detector quality ablation ($\rho=0.4$, seed42). A 5-point AUC improvement yields only 0.005 PPL gain.}
\label{tab:detector}
\begin{tabular}{lcccc}
\toprule
\textbf{Detector} & \textbf{AUC} & \textbf{Adaptive PPL} & $\boldsymbol{\Delta}$ \textbf{vs.\ unif.} \\
\midrule
TF-IDF + LR & 0.939 & 17.650 & $-$0.040 \\
ModernBERT & 0.986 & 17.645 & $-$0.045 \\
\midrule
$\Delta_{\text{detector}}$ & +0.047 & $-$0.005 & --- \\
\bottomrule
\end{tabular}
\end{table}

Table~\ref{tab:detector} reveals that a 5-point AUC gain (0.939\,$\to$\,0.986) from replacing TF-IDF+LR with ModernBERT yields only 0.005 PPL improvement---far below seed variance.
The marginal return is minimal because adaptive reweighting depends on $q$-score \emph{ranking} rather than absolute accuracy: as long as rank ordering is approximately correct (AUC\,$>$\,0.9), $\bstar$ derivation functions effectively.
This finding has practical implications: a lightweight detector running in minutes suffices, eliminating the need for expensive neural detectors.

\subsection{Cross-Generator Robustness}
\label{sec:cross-gen}

\begin{table}[t]
\centering
\small
\caption{In-distribution vs.\ cross-generator contamination ($\rho=0.4$, seed42).}
\label{tab:cross-gen}
\begin{tabular}{lccccc}
\toprule
\textbf{Generator} & \textbf{Det.\ AUC} & \textbf{Unif.} & \textbf{Adapt.} & $\boldsymbol{\Delta}$ \\
\midrule
Qwen3-0.6B (in-dist.) & 0.939 & 17.69 & 17.65 & $-$0.04 \\
Llama-3.2-1B (cross) & 0.925 & 17.72 & 17.69 & $-$0.03 \\
\bottomrule
\end{tabular}
\end{table}

Cross-generator contamination reduces adaptive improvement from 0.04 to 0.03 PPL (Table~\ref{tab:cross-gen}).
The reduction stems from the detector's slightly lower AUC on out-of-distribution synthetic text (0.925 vs.\ 0.939) and increased $q$-score distribution dispersion ($\sigma$: 0.37\,$\to$\,0.39), which degrades per-batch $\bstar$ estimation.
The method remains directionally effective but with reduced magnitude.

\subsection{Domain-Skewed Contamination}
\label{sec:domain-skew}

When contamination is concentrated in a single domain (code-heavy injection), adaptive improvement drops to 0.02 PPL (vs.\ 0.04 for uniform injection), achieving only $\sim$50\% of the uniform-injection ratio.
Code-heavy batches exhibit near-uniform high $q$-scores ($\sigma \approx 0.05$), violating the implicit assumption that batches contain sufficient $q$-score variance for meaningful $\bstar$ differentiation.
This represents a limitation of per-batch $\bstar$ estimation (\S\ref{sec:limitations}).

\subsection{Downstream Evaluation}
\label{sec:downstream}

\begin{table}[t]
\centering
\small
\caption{Downstream tasks (Llama-3.2-3B, $\rho=0.4$, seed42, 200M tokens).}
\label{tab:downstream}
\begin{tabular}{lcccc}
\toprule
\textbf{Task} & \textbf{Uniform} & \textbf{Adaptive} & $\boldsymbol{\Delta}$ & \textbf{Sig.} \\
\midrule
MMLU (5-shot) & 34.20 & 34.35 & +0.15 & No \\
HellaSwag (0-shot) & 52.80 & 53.10 & +0.30 & Marginal \\
eval PPL & 13.52 & 13.46 & $-$0.06 & Yes \\
\bottomrule
\end{tabular}
\end{table}

Table~\ref{tab:downstream} shows downstream evaluation on MMLU and HellaSwag.
Improvements are directionally positive but not statistically significant.
At 200M training tokens, the 3B model is far from convergence on these benchmarks (MMLU baseline 34.2\% vs.\ 25\% random), and the 0.44\% relative PPL improvement is insufficient to produce detectable downstream signal.
We expect downstream separation to emerge at larger token budgets ($>$1B), which we leave to future work.

% ============================================================
\section{Analysis}
\label{sec:analysis}

\paragraph{Binding constraint statistics.}
Across 51,913 micro-batches (seed42, $\rho=0.4$), 49.9\% are ``clean'' ($\bstar=0$), 20.2\% have $\aeff$ as the binding constraint, and 30.0\% have ESS as binding.
Among reweighted batches, $\aeff$ binds 40.1\% of the time with mean $\bstar=1.67$, while ESS-only would select $\bstar=4.15$ (2.5$\times$ more aggressive).
At $\aeff$-binding points, ESS averages 2.60 (vs.\ floor $\nmin=2.0$), retaining 30\% headroom---demonstrating that the minimum-intervention principle preserves substantial effective sample diversity relative to max-intervention.

\paragraph{Flat PPL landscape.}
The extended grid reveals that PPL varies by only 0.014 across $b\in[0.1,0.5]$.
This flatness explains why dual-constraint $\approx$ ESS-only empirically: the two strategies select different $\bstar$ values (median 0.88 vs.\ 3.20) that happen to lie on the same plateau.
We recommend dual-constraint as the default not for PPL advantage but for lower gradient variance and principled theoretical grounding.

\paragraph{$\bstar$ adapts to contamination.}
At $\rho=0.6$, $\bstar$ median shifts from 0.88 to 0.62 as higher contamination causes ESS to bind more frequently (72\% vs.\ 60\%).
This automatic adaptation---without any hyperparameter change---is a direct consequence of computing $\bstar$ from batch statistics.

% ============================================================
\section{Related Work}
\label{sec:related}

\paragraph{Synthetic data contamination.}
\citet{strong-collapse-iclr2025} establish that recursive training on model-generated text leads to model collapse even with small synthetic fractions, motivating contamination-aware training strategies.

\paragraph{Detect-and-reweight.}
\citet{drayson-emnlp2025} propose the $w=(1-q)^b$ reweighting scheme with SIR resampling, demonstrating effectiveness on WikiText-2 at 2M tokens.
Their approach requires per-model validation to select $b$.
We adopt the same weighting function but derive $b$ adaptively via dual constraints and apply it through loss reweighting at 100$\times$ larger scale.

\paragraph{Data mixing and weighting.}
DoReMi~\citep{doremi} and DSIR~\citep{dsir} optimize \emph{domain}-level mixing weights, while we operate at the \emph{instance} level.
SoftDedup~\citep{softdedup} applies repetition-based soft reweighting; our approach is orthogonal, targeting provenance rather than duplication.
Dynamic loss-based reweighting~\citep{dynamic-loss} uses training loss itself as the weighting signal, without external detector scores.

\paragraph{Importance sampling.}
Kish's effective sample size and Pareto-smoothed importance sampling~\citep{psis} provide the theoretical foundation for our ESS constraint.
We apply these diagnostics \emph{prescriptively} (as an optimization constraint) rather than descriptively.

% ============================================================
\section{Limitations}
\label{sec:limitations}

\paragraph{Scale.}
Experiments use 200M tokens and models up to 3B parameters.
While we observe consistent trends across model scales and contamination ratios, validation at $>$1B tokens and $>$7B parameters is needed.

\paragraph{Per-batch estimation.}
The $\bstar$ derivation assumes sufficient $q$-score variance within each micro-batch.
Domain-skewed contamination (e.g., code-heavy injection) produces batches with near-uniform $q$-scores, reducing $\bstar$ effectiveness to $\sim$50\% of the uniform-injection case.
Larger batch sizes or cross-batch smoothing may mitigate this.

\paragraph{Downstream transfer.}
PPL improvements do not yet translate to statistically significant downstream gains at our training scale.

\paragraph{$\tau$ boundaries.}
While $\tau\in[0.6,0.8]$ is robust, boundary values ($\tau=0.5$ and $\tau\geq 0.9$) degrade by 0.03--0.07 PPL.
The method is not fully $\tau$-insensitive.

\paragraph{Contamination regime.}
At $\rho=0.6$, adaptive improvement amplifies only modestly (1.25$\times$), with ESS constraints limiting additional headroom.
The method mitigates but does not solve high-contamination degradation (18.45 vs.\ 18.50 uniform, on a baseline of 17.69 at $\rho=0.4$).

% ============================================================
\section{Conclusion}
\label{sec:conclusion}

We presented dual-constrained adaptive $\bstar$, a principled framework for selecting the reweighting exponent in contamination-aware pretraining.
By jointly enforcing a quality lower bound ($\aeff\geq\tau$) and a stability upper bound ($\mathrm{ESS}\geq\nmin$), the method implements a minimum-intervention principle that requires no validation-based grid search and transfers across model scales without tuning.
On controlled experiments spanning two models (0.6B, 3B), two contamination ratios (40\%, 60\%), and multiple batch sizes, adaptive $\bstar$ matches grid-search-optimal performance within seed variance.
A lightweight TF-IDF+LR detector (AUC\,=\,0.939) suffices---the marginal return from stronger detection is negligible.
The adaptive $\bstar$ framework, including its ESS-only variant, eliminates the per-model grid search cost entirely, providing a practical and principled approach to training on contaminated web corpora.

% ============================================================
\begin{thebibliography}{10}
\providecommand{\natexlab}[1]{#1}

\bibitem[Abbas et~al.(2024)]{softdedup}
Abbas, A., Tirumala, K., Simig, D., Ganguli, S., and Morcos, A.~S. (2024).
\newblock {SemDeDup} meets {SoftDeDup}: Data-efficient training via soft deduplication.
\newblock In \emph{Proc.\ ACL}.

\bibitem[Chen et~al.(2025)]{convergence-theory}
Chen, Y., Fan, J., and Ma, C. (2025).
\newblock Convergence rates for contaminated mixture estimation.
\newblock \emph{arXiv:2602.16065}.

\bibitem[Drayson et~al.(2025)]{drayson-emnlp2025}
Drayson, R., Wenger, E., and Tram\`{e}r, F. (2025).
\newblock Detect-and-reweight: Provenance-based training data curation for language models.
\newblock In \emph{Proc.\ EMNLP}.

\bibitem[Fan et~al.(2025)]{dynamic-loss}
Fan, Y., Chen, Y., and Du, S. (2025).
\newblock Dynamic loss-based data reweighting for improved language model pretraining.
\newblock In \emph{Proc.\ ICLR}.

\bibitem[Li et~al.(2025)]{golden-ratio}
Li, Y., Pan, J., and Gao, Y. (2025).
\newblock Golden ratio weighting for mixed real-synthetic pretraining.
\newblock \emph{arXiv:2502.18049}.

\bibitem[Penedo et~al.(2024)]{fineweb}
Penedo, G., Kydl\'{i}\v{c}ek, H., Ben~Allal, L., Lozhkov, A., Mitchell, M., Raffel, C., von Werra, L., and Wolf, T. (2024).
\newblock The {FineWeb} datasets: Decanting the web for the finest text data at scale.
\newblock In \emph{Proc.\ NeurIPS Datasets and Benchmarks Track}.

\bibitem[Shumailov et~al.(2025)]{strong-collapse-iclr2025}
Shumailov, I., Shumaylov, Z., Zhao, Y., Papernot, N., Anderson, R., and Gal, Y. (2025).
\newblock {AI} models collapse when trained on recursively generated data.
\newblock In \emph{Proc.\ ICLR}.

\bibitem[Vehtari et~al.(2024)]{psis}
Vehtari, A., Gelman, A., Simpson, D., Carpenter, B., and B\"{u}rkner, P.-C. (2024).
\newblock Pareto smoothed importance sampling.
\newblock \emph{Journal of Machine Learning Research}, 25(72):1--58.

\bibitem[Xie et~al.(2023a)]{doremi}
Xie, S.~M., Pham, H., Dong, X., Du, N., Liu, H., Lu, Y., Liang, P., Le, Q.~V., Ma, T., and Yu, A.~W. (2023a).
\newblock {DoReMi}: Optimizing data mixtures speeds up language model pretraining.
\newblock In \emph{Proc.\ NeurIPS}.

\bibitem[Xie et~al.(2023b)]{dsir}
Xie, S.~M., Santurkar, S., Ma, T., and Liang, P. (2023b).
\newblock Data selection for language models via importance resampling.
\newblock In \emph{Proc.\ NeurIPS}.

\end{thebibliography}

\end{document}
